%%
% 摘要信息
% 摘要内容应概括地反映出本论文的主要内容，主要说明本论文的研究目的、内容、方法、成果和结论。要突出本论文的创造性成果或新见解，不要与引言相 混淆。语言力求精练、准确，以 300—500 字为宜。
% 关键词是供检索用的主题词条，应采用能覆盖论文主要内容的通用技术词条(参照相应的技术术语 标准)。按词条的外延层次排列（外延大的排在前面）。


\cabstract{
随着B5G/6G网络向海量终端连接与多样化业务需求持续演进，复杂物联网场景下随机接入过程中的性能瓶颈日益凸显。传统Aloha与CSMA/LBT等接入机制普遍存在按包重复竞争开销较大、信道利用率受限以及对动态业务变化适应能力不足等问题，严重制约系统在高负载与多业务异构环境中的整体性能表现。围绕复杂物联网场景下高效、低开销且可自适应的随机接入核心需求，本文从接入机制优化与系统级决策两个维度开展系统性研究，为B5G/6G网络随机接入系统的性能提升与工程实现提供理论支撑与技术路径。

面向授权频谱短报文接入场景，针对传统Aloha机制重复竞争开销偏高、适配性不足的问题，提出连续传输增强Aloha机制。该机制在保持原有分布式接入结构不变的前提下，通过扩展单次竞争成功后的连续发送能力，有效降低单位数据传输的平均竞争开销，兼顾机制兼容性与性能提升需求。基于休假排队理论构建系统性能分析模型，对吞吐量、信令开销等关键指标进行量化刻画，并结合离散事件仿真验证机制的有效性与可行性。研究结果表明，所提机制可将系统最大吞吐量提升至0.5，突破传统机制性能上限，同时在5G小数据传输场景中显著降低信令开销，具备良好的工程兼容性与实际部署价值。

面向非授权频谱LBT接入场景，针对传统CSMA/LBT机制因频繁侦听、退避导致的系统开销过大、性能受限问题，提出连续传输增强CSMA机制。该机制在完整保留载波侦听与退避核心流程的基础上，引入连续传输策略，有效降低重复竞争带来的资源消耗与效率损失。针对信道侦听与传输时长异构的核心问题，构建基于吸收马尔可夫链的精细化分析模型，对系统吞吐量与接入公平性进行联合刻画与量化分析。理论分析与仿真验证结果表明，该机制在无需额外复杂控制逻辑的前提下，可实现约28.8\%的吞吐量提升，且在效率提升的同时稳定保持良好的系统公平性，适配非授权频谱复杂接入场景需求。

针对动态业务环境中单一接入机制难以持续适配场景变化、无法兼顾多维度性能需求的问题，构建统一的接入性能评估与自适应决策框架。通过离散事件仿真技术，生成覆盖多协议、多参数配置的大规模性能数据集，基于多任务学习模型实现吞吐量、时延、公平性及信令开销等多维性能指标的快速精准预测。在此基础上，设计面向多目标优化的效用函数，实现接入协议选择与参数配置的协同优化，提出自适应决策方法。实验结果表明，该方法在保证较低计算开销的前提下，能够实现动态场景下的快速决策与稳定性能增益，显著提升系统对动态业务环境的适配能力，具备较强的实际部署潜力。

综上，本文从机制设计与系统决策两个层面出发，构建了一套面向复杂物联网场景的低开销、高效率且具备自适应能力的随机接入解决方案，系统提升了B5G/6G 下随机接入的系统综合性能，为随机接入增强机制的工程实现与系统优化提供了重要的理论参考与技术支撑。
}
% 中文关键词(每个关键词之间用;分开,最后一个关键词不打标点符号。)
\ckeywords{5G；随机接入；Aloha；CSMA/LBT；连续传输；休假排队；多任务学习；自适应决策}

\eabstract{
With the continuous evolution of B5G/6G networks towards massive terminal connections and diverse service requirements, the performance bottlenecks in the random access process of complex Internet of Things scenarios have become increasingly prominent. Traditional access mechanisms such as Aloha and CSMA/LBT generally have problems including high overhead caused by packet-by-packet repeated contention, limited channel utilization, and insufficient adaptability to dynamic service changes, which severely restrict the overall performance of the system in high-load and multi-service heterogeneous environments. Centering on the core demand for efficient, low-overhead and adaptive random access in complex Internet of Things scenarios, this paper conducts systematic research from two dimensions: access mechanism optimization and system-level decision-making, providing theoretical support and technical paths for the performance improvement and engineering implementation of random access systems in B5G/6G networks.

For the short packet access scenario in authorized spectrum, aiming at the problems of high repeated contention overhead and insufficient adaptability of traditional Aloha mechanisms, a successive transmission enhanced Aloha mechanism is proposed. On the premise of maintaining the original distributed access structure unchanged, the mechanism effectively reduces the average contention cost of unit data transmission by extending the continuous transmission capability after a single successful contention, balancing mechanism compatibility and performance improvement requirements. A system performance analysis model is constructed based on vacation queuing theory to quantitatively characterize key indicators such as throughput and signaling overhead, and the effectiveness and feasibility of the mechanism are verified through discrete event simulation. The research results show that the proposed mechanism can increase the maximum system throughput to 0.5, break the performance upper limit of traditional mechanisms, and significantly reduce signaling overhead in 5G small data transmission scenarios, with good engineering compatibility and practical deployment value.

For the LBT access scenario in unlicensed spectrum, aiming at the problems of excessive system overhead and limited performance of traditional CSMA/LBT mechanisms caused by frequent listening and backoff, a successive transmission enhanced CSMA mechanism is proposed. On the basis of fully retaining the core processes of carrier sensing and backoff, the mechanism introduces a successive transmission strategy to effectively reduce resource consumption and efficiency loss caused by repeated contention. Aiming at the core problem of heterogeneous channel listening and transmission durations, a refined analysis model based on absorbing Markov chain is constructed to jointly characterize and quantitatively analyze system throughput and access fairness. Theoretical analysis and simulation verification results show that without additional complex control logic, the proposed mechanism can achieve a throughput improvement of about 28.8\%, and stably maintain good system fairness while improving efficiency, adapting to the needs of complex access scenarios in unlicensed spectrum.

Aiming at the problem that a single access mechanism is difficult to continuously adapt to scenario changes and cannot balance multi-dimensional performance requirements in dynamic business environments, a unified access performance evaluation and adaptive decision-making framework is constructed. Large-scale performance datasets covering multi-protocol and multi-parameter configurations are generated through discrete event simulation technology, and a multi-task learning model is used to achieve fast and accurate prediction of multi-dimensional performance indicators including throughput, delay, fairness and signaling overhead. On this basis, a utility function oriented to multi-objective optimization is designed to realize the coordinated optimization of access protocol selection and parameter configuration, and an adaptive decision-making method is proposed. Experimental results show that on the premise of ensuring low computational overhead, the method can achieve fast decision-making and stable performance gains in dynamic scenarios, significantly improve the adaptability of the system to dynamic business environments, and has strong practical deployment potential.

In summary, starting from the two levels of mechanism design and system decision-making, this paper constructs a set of low-overhead, high-efficiency and adaptive random access solutions for complex Internet of Things scenarios, systematically improving the comprehensive performance of random access systems in B5G/6G networks, and providing important theoretical reference and technical support for the engineering implementation and system optimization of random access mechanisms.
}
% 英文文关键词（关键词之间用逗号隔开，最后一个关键词不打标点符号。）
\ekeywords{ 5G; Random Access, Aloha, CSMA/LBT, Successive Transmission, Vacation Queuing Theory, Multi-Task Learning, Adaptive Decision-Making} 
